% =============================================================================
% main.tex -- Entry point Template Tesis Magister IPB University
%
% Template LaTeX ini diadaptasi dari PPTA IPB 2026 dan templat resmi
% "format tesis terbaru" untuk program magister, dengan basis struktur
% ipb-thesis-latex karya Raihan Putra Kirana (@raihanpka).
%
% Struktur:
%   src/
%     main.tex                    Berkas ini (entry point)
%     config/                     Paket dan konfigurasi format
%       ipb-thesis.sty            Style utama (margins, font, headings, dll.)
%       information.tex           Data penulis, tesis, institusi
%       hyphenation-id.tex        Aturan pemenggalan kata
%     ringkasan/                  Ringkasan (id) dan Summary (en)
%     preliminaries/              Bagian Awal (cover, pernyataan, dll.)
%     chapters/                   Bagian Inti (Bab I s.d. Bab V)
%     backmatter/                 Bagian Akhir (daftar pustaka, lampiran)
%     resources/                  Gambar, logo, dan sumber daya lain
%     refs/                       Database referensi BibTeX
%
% Lisensi: MIT (lihat LICENSE). Kredit lengkap ada di README.md.
% =============================================================================

% -----------------------------------------------------------------------------
% 1. Kelas Dokumen & Paket Utama
% -----------------------------------------------------------------------------
\documentclass[12pt,a4paper,twoside,openright]{book}

% Style utama (margins, font, headings, captions, dst.)
\usepackage{ipb-thesis}

% Data penulis, tesis, institusi
% =============================================================================
% information.tex
%
% Satu sumber kebenaran untuk data penulis, tesis, institusi, pernyataan AI,
% pengesahan, dan komponen awal lain sesuai PPTA IPB 2026 dan templat
% "format tesis terbaru" (program magister).
%
% Pemuatan: % =============================================================================
% information.tex
%
% Satu sumber kebenaran untuk data penulis, tesis, institusi, pernyataan AI,
% pengesahan, dan komponen awal lain sesuai PPTA IPB 2026 dan templat
% "format tesis terbaru" (program magister).
%
% Pemuatan: % =============================================================================
% information.tex
%
% Satu sumber kebenaran untuk data penulis, tesis, institusi, pernyataan AI,
% pengesahan, dan komponen awal lain sesuai PPTA IPB 2026 dan templat
% "format tesis terbaru" (program magister).
%
% Pemuatan: \input{config/information} dari src/main.tex.
% =============================================================================

% --- Opsi Tampilan ---
% Catatan panduan berwarna merah berasal dari templat Word resmi: teks
% instruksi yang harus dihapus pada naskah akhir. Ubah satu baris berikut
% menjadi false agar semua catatan merah hilang sekaligus (judul Daftar
% Tabel/Gambar/Lampiran, Prakata, dan catatan lain yang memakai \ipbcatatan).
% Tata letak tetap sama karena ruang catatan dipertahankan.
\renewcommand{\TampilkanCatatanPanduan}{true}

% --- Data Penulis ---
% Nilai bawaan mengikuti teks placeholder pada templat resmi. Ganti seluruh
% placeholder ini dengan data sebenarnya sebelum naskah diserahkan.
\newcommand{\NamaPenulis}{Nama Penulis}
\newcommand{\NamaPenulisInggris}{Student Name}
\newcommand{\NIM}{NIM}
\newcommand{\TempatLahir}{Kota \ldots{}}
\newcommand{\TanggalLahir}{tanggal bulan tahun}
\newcommand{\EmailPenulis}{alamat.surel@apps.ipb.ac.id}

% --- Data Pembimbing ---
\newcommand{\NamaPembimbingSatu}{Nama Pembimbing 1}
\newcommand{\NIPPembimbingSatu}{}
\newcommand{\NamaPembimbingDua}{Nama Pembimbing 2}
\newcommand{\NIPPembimbingDua}{}
% Pembimbing ketiga bersifat opsional. Set \TampilkanPembimbingTiga menjadi
% true bila komisi pembimbing beranggotakan tiga orang. Nilai ini mengatur
% baris tanda tangan pada halaman pengesahan sekaligus penyebutan nama pada
% Sorotan/Highlights, Ringkasan/Summary, dan Prakata.
\newcommand{\TampilkanPembimbingTiga}{false}
\newcommand{\NamaPembimbingTiga}{Nama Pembimbing 3}

% Daftar pembimbing yang dipakai di dalam kalimat.
\newcommand{\DaftarPembimbing}{%
  \ifdefstring{\TampilkanPembimbingTiga}{true}%
    {\NamaPembimbingSatu, \NamaPembimbingDua, dan \NamaPembimbingTiga}%
    {\NamaPembimbingSatu{} dan \NamaPembimbingDua}%
}
% Versi kapital untuk baris nama pada Sorotan/Highlights dan Ringkasan/Summary.
\newcommand{\DaftarPembimbingKapital}{%
  \ifdefstring{\TampilkanPembimbingTiga}{true}%
    {\MakeUppercase{\NamaPembimbingSatu}, \MakeUppercase{\NamaPembimbingDua},
     dan \MakeUppercase{\NamaPembimbingTiga}}%
    {\MakeUppercase{\NamaPembimbingSatu} dan \MakeUppercase{\NamaPembimbingDua}}%
}
% Versi bahasa Inggris ("and") untuk baris nama pada Highlights/Summary.
\newcommand{\DaftarPembimbingKapitalInggris}{%
  \ifdefstring{\TampilkanPembimbingTiga}{true}%
    {\MakeUppercase{\NamaPembimbingSatu}, \MakeUppercase{\NamaPembimbingDua},
     and \MakeUppercase{\NamaPembimbingTiga}}%
    {\MakeUppercase{\NamaPembimbingSatu} and \MakeUppercase{\NamaPembimbingDua}}%
}

% --- Data Penguji ---
\newcommand{\NamaPengujiSatu}{Nama Lengkap dan Gelar}
\newcommand{\NamaPengujiDua}{Nama Lengkap dan Gelar}

% --- Data Pengesahan ---
\newcommand{\NamaKetuaProgramStudi}{Nama Lengkap dan Gelar}
\newcommand{\NIPKetuaProgramStudi}{}
% Pejabat kedua pada halaman pengesahan: Dekan Fakultas/Sekolah.
\newcommand{\JabatanPejabatKedua}{Dekan Fakultas/Sekolah}
\newcommand{\NamaPejabatKedua}{Nama Lengkap dan Gelar}
\newcommand{\NIPPejabatKedua}{}
% Biarkan kosong agar templat hanya menampilkan label dan ruang pengisian.
\newcommand{\TanggalUjian}{}
\newcommand{\TanggalLulus}{}

% --- Data Tesis ---
\newcommand{\JudulTesis}{%
  Judul Karya Ilmiah Maksimum Tiga Baris, Lima Belas Kata Tidak Termasuk
  Kata Depan dan Kata Sambung%
}
\newcommand{\JudulTesisInggris}{%
  Title of Master's Thesis%
}
\newcommand{\KataKunci}{ditulis dalam bahasa Indonesia, disusun berdasarkan abjad, maksimum lima kata atau frasa}
\newcommand{\Keywords}{written in English, arranged alphabetically, maximum five words or phrases}

% Judul ringkasan visual maksimum enam kata, tidak termasuk kata depan/sambung.
\newcommand{\JudulRingkasanVisual}{%
  Judul Ringkasan Visual Maksimum Enam Kata, Tidak Termasuk Kata Depan dan Kata Sambung}
\newcommand{\JudulRingkasanVisualInggris}{%
  Graphical Summary Title Maximum Six Words, Excluding Prepositions and Conjunctions}
% Isi dengan path berkas gambar untuk mengganti kotak placeholder.
\newcommand{\BerkasRingkasanVisual}{}
\newcommand{\BerkasRingkasanVisualInggris}{}

% --- Deklarasi penggunaan AI ---
% Nilai placeholder mempertahankan pilihan [menggunakan/tidak menggunakan]
% seperti pada DOTX. Ganti dengan true atau false setelah menentukan pilihan.
\newcommand{\MenggunakanAI}{placeholder}
\newcommand{\NamaAlatAI}{[nama alat/layanan]}
\newcommand{\TujuanPenggunaanAI}{[alasan]}

% --- Gelar dan sumber dana ---
\newcommand{\GelarMagister}{magister \ldots{}}
% Kosongkan bila penelitian tidak memperoleh pendanaan khusus.
\newcommand{\SumberDana}{}

% --- Data Institusi ---
\newcommand{\Departemen}{Nama Departemen}
\newcommand{\ProgramStudi}{Nama Program Studi}
\newcommand{\ProgramStudiSingkat}{Nama Program Studi}
\newcommand{\Fakultas}{Fakultas/Sekolah}
\newcommand{\Universitas}{Institut Pertanian Bogor}
\newcommand{\KotaUniversitas}{Bogor}
\newcommand{\Tahun}{20XX}

% --- Data Waktu Pelaksanaan (untuk prakata) ---
\newcommand{\BulanMulaiPenelitian}{\ldots{}}
\newcommand{\TahunMulaiPenelitian}{20XX}
\newcommand{\BulanSelesaiPenelitian}{\ldots{}}
\newcommand{\TahunSelesaiPenelitian}{20XX}
\newcommand{\BulanTahunDokumen}{Bulan Tahun}
 dari src/main.tex.
% =============================================================================

% --- Opsi Tampilan ---
% Catatan panduan berwarna merah berasal dari templat Word resmi: teks
% instruksi yang harus dihapus pada naskah akhir. Ubah satu baris berikut
% menjadi false agar semua catatan merah hilang sekaligus (judul Daftar
% Tabel/Gambar/Lampiran, Prakata, dan catatan lain yang memakai \ipbcatatan).
% Tata letak tetap sama karena ruang catatan dipertahankan.
\renewcommand{\TampilkanCatatanPanduan}{true}

% --- Data Penulis ---
% Nilai bawaan mengikuti teks placeholder pada templat resmi. Ganti seluruh
% placeholder ini dengan data sebenarnya sebelum naskah diserahkan.
\newcommand{\NamaPenulis}{Nama Penulis}
\newcommand{\NamaPenulisInggris}{Student Name}
\newcommand{\NIM}{NIM}
\newcommand{\TempatLahir}{Kota \ldots{}}
\newcommand{\TanggalLahir}{tanggal bulan tahun}
\newcommand{\EmailPenulis}{alamat.surel@apps.ipb.ac.id}

% --- Data Pembimbing ---
\newcommand{\NamaPembimbingSatu}{Nama Pembimbing 1}
\newcommand{\NIPPembimbingSatu}{}
\newcommand{\NamaPembimbingDua}{Nama Pembimbing 2}
\newcommand{\NIPPembimbingDua}{}
% Pembimbing ketiga bersifat opsional. Set \TampilkanPembimbingTiga menjadi
% true bila komisi pembimbing beranggotakan tiga orang. Nilai ini mengatur
% baris tanda tangan pada halaman pengesahan sekaligus penyebutan nama pada
% Sorotan/Highlights, Ringkasan/Summary, dan Prakata.
\newcommand{\TampilkanPembimbingTiga}{false}
\newcommand{\NamaPembimbingTiga}{Nama Pembimbing 3}

% Daftar pembimbing yang dipakai di dalam kalimat.
\newcommand{\DaftarPembimbing}{%
  \ifdefstring{\TampilkanPembimbingTiga}{true}%
    {\NamaPembimbingSatu, \NamaPembimbingDua, dan \NamaPembimbingTiga}%
    {\NamaPembimbingSatu{} dan \NamaPembimbingDua}%
}
% Versi kapital untuk baris nama pada Sorotan/Highlights dan Ringkasan/Summary.
\newcommand{\DaftarPembimbingKapital}{%
  \ifdefstring{\TampilkanPembimbingTiga}{true}%
    {\MakeUppercase{\NamaPembimbingSatu}, \MakeUppercase{\NamaPembimbingDua},
     dan \MakeUppercase{\NamaPembimbingTiga}}%
    {\MakeUppercase{\NamaPembimbingSatu} dan \MakeUppercase{\NamaPembimbingDua}}%
}
% Versi bahasa Inggris ("and") untuk baris nama pada Highlights/Summary.
\newcommand{\DaftarPembimbingKapitalInggris}{%
  \ifdefstring{\TampilkanPembimbingTiga}{true}%
    {\MakeUppercase{\NamaPembimbingSatu}, \MakeUppercase{\NamaPembimbingDua},
     and \MakeUppercase{\NamaPembimbingTiga}}%
    {\MakeUppercase{\NamaPembimbingSatu} and \MakeUppercase{\NamaPembimbingDua}}%
}

% --- Data Penguji ---
\newcommand{\NamaPengujiSatu}{Nama Lengkap dan Gelar}
\newcommand{\NamaPengujiDua}{Nama Lengkap dan Gelar}

% --- Data Pengesahan ---
\newcommand{\NamaKetuaProgramStudi}{Nama Lengkap dan Gelar}
\newcommand{\NIPKetuaProgramStudi}{}
% Pejabat kedua pada halaman pengesahan: Dekan Fakultas/Sekolah.
\newcommand{\JabatanPejabatKedua}{Dekan Fakultas/Sekolah}
\newcommand{\NamaPejabatKedua}{Nama Lengkap dan Gelar}
\newcommand{\NIPPejabatKedua}{}
% Biarkan kosong agar templat hanya menampilkan label dan ruang pengisian.
\newcommand{\TanggalUjian}{}
\newcommand{\TanggalLulus}{}

% --- Data Tesis ---
\newcommand{\JudulTesis}{%
  Judul Karya Ilmiah Maksimum Tiga Baris, Lima Belas Kata Tidak Termasuk
  Kata Depan dan Kata Sambung%
}
\newcommand{\JudulTesisInggris}{%
  Title of Master's Thesis%
}
\newcommand{\KataKunci}{ditulis dalam bahasa Indonesia, disusun berdasarkan abjad, maksimum lima kata atau frasa}
\newcommand{\Keywords}{written in English, arranged alphabetically, maximum five words or phrases}

% Judul ringkasan visual maksimum enam kata, tidak termasuk kata depan/sambung.
\newcommand{\JudulRingkasanVisual}{%
  Judul Ringkasan Visual Maksimum Enam Kata, Tidak Termasuk Kata Depan dan Kata Sambung}
\newcommand{\JudulRingkasanVisualInggris}{%
  Graphical Summary Title Maximum Six Words, Excluding Prepositions and Conjunctions}
% Isi dengan path berkas gambar untuk mengganti kotak placeholder.
\newcommand{\BerkasRingkasanVisual}{}
\newcommand{\BerkasRingkasanVisualInggris}{}

% --- Deklarasi penggunaan AI ---
% Nilai placeholder mempertahankan pilihan [menggunakan/tidak menggunakan]
% seperti pada DOTX. Ganti dengan true atau false setelah menentukan pilihan.
\newcommand{\MenggunakanAI}{placeholder}
\newcommand{\NamaAlatAI}{[nama alat/layanan]}
\newcommand{\TujuanPenggunaanAI}{[alasan]}

% --- Gelar dan sumber dana ---
\newcommand{\GelarMagister}{magister \ldots{}}
% Kosongkan bila penelitian tidak memperoleh pendanaan khusus.
\newcommand{\SumberDana}{}

% --- Data Institusi ---
\newcommand{\Departemen}{Nama Departemen}
\newcommand{\ProgramStudi}{Nama Program Studi}
\newcommand{\ProgramStudiSingkat}{Nama Program Studi}
\newcommand{\Fakultas}{Fakultas/Sekolah}
\newcommand{\Universitas}{Institut Pertanian Bogor}
\newcommand{\KotaUniversitas}{Bogor}
\newcommand{\Tahun}{20XX}

% --- Data Waktu Pelaksanaan (untuk prakata) ---
\newcommand{\BulanMulaiPenelitian}{\ldots{}}
\newcommand{\TahunMulaiPenelitian}{20XX}
\newcommand{\BulanSelesaiPenelitian}{\ldots{}}
\newcommand{\TahunSelesaiPenelitian}{20XX}
\newcommand{\BulanTahunDokumen}{Bulan Tahun}
 dari src/main.tex.
% =============================================================================

% --- Opsi Tampilan ---
% Catatan panduan berwarna merah berasal dari templat Word resmi: teks
% instruksi yang harus dihapus pada naskah akhir. Ubah satu baris berikut
% menjadi false agar semua catatan merah hilang sekaligus (judul Daftar
% Tabel/Gambar/Lampiran, Prakata, dan catatan lain yang memakai \ipbcatatan).
% Tata letak tetap sama karena ruang catatan dipertahankan.
\renewcommand{\TampilkanCatatanPanduan}{true}

% --- Data Penulis ---
% Nilai bawaan mengikuti teks placeholder pada templat resmi. Ganti seluruh
% placeholder ini dengan data sebenarnya sebelum naskah diserahkan.
\newcommand{\NamaPenulis}{Nama Penulis}
\newcommand{\NamaPenulisInggris}{Student Name}
\newcommand{\NIM}{NIM}
\newcommand{\TempatLahir}{Kota \ldots{}}
\newcommand{\TanggalLahir}{tanggal bulan tahun}
\newcommand{\EmailPenulis}{alamat.surel@apps.ipb.ac.id}

% --- Data Pembimbing ---
\newcommand{\NamaPembimbingSatu}{Nama Pembimbing 1}
\newcommand{\NIPPembimbingSatu}{}
\newcommand{\NamaPembimbingDua}{Nama Pembimbing 2}
\newcommand{\NIPPembimbingDua}{}
% Pembimbing ketiga bersifat opsional. Set \TampilkanPembimbingTiga menjadi
% true bila komisi pembimbing beranggotakan tiga orang. Nilai ini mengatur
% baris tanda tangan pada halaman pengesahan sekaligus penyebutan nama pada
% Sorotan/Highlights, Ringkasan/Summary, dan Prakata.
\newcommand{\TampilkanPembimbingTiga}{false}
\newcommand{\NamaPembimbingTiga}{Nama Pembimbing 3}

% Daftar pembimbing yang dipakai di dalam kalimat.
\newcommand{\DaftarPembimbing}{%
  \ifdefstring{\TampilkanPembimbingTiga}{true}%
    {\NamaPembimbingSatu, \NamaPembimbingDua, dan \NamaPembimbingTiga}%
    {\NamaPembimbingSatu{} dan \NamaPembimbingDua}%
}
% Versi kapital untuk baris nama pada Sorotan/Highlights dan Ringkasan/Summary.
\newcommand{\DaftarPembimbingKapital}{%
  \ifdefstring{\TampilkanPembimbingTiga}{true}%
    {\MakeUppercase{\NamaPembimbingSatu}, \MakeUppercase{\NamaPembimbingDua},
     dan \MakeUppercase{\NamaPembimbingTiga}}%
    {\MakeUppercase{\NamaPembimbingSatu} dan \MakeUppercase{\NamaPembimbingDua}}%
}
% Versi bahasa Inggris ("and") untuk baris nama pada Highlights/Summary.
\newcommand{\DaftarPembimbingKapitalInggris}{%
  \ifdefstring{\TampilkanPembimbingTiga}{true}%
    {\MakeUppercase{\NamaPembimbingSatu}, \MakeUppercase{\NamaPembimbingDua},
     and \MakeUppercase{\NamaPembimbingTiga}}%
    {\MakeUppercase{\NamaPembimbingSatu} and \MakeUppercase{\NamaPembimbingDua}}%
}

% --- Data Penguji ---
\newcommand{\NamaPengujiSatu}{Nama Lengkap dan Gelar}
\newcommand{\NamaPengujiDua}{Nama Lengkap dan Gelar}

% --- Data Pengesahan ---
\newcommand{\NamaKetuaProgramStudi}{Nama Lengkap dan Gelar}
\newcommand{\NIPKetuaProgramStudi}{}
% Pejabat kedua pada halaman pengesahan: Dekan Fakultas/Sekolah.
\newcommand{\JabatanPejabatKedua}{Dekan Fakultas/Sekolah}
\newcommand{\NamaPejabatKedua}{Nama Lengkap dan Gelar}
\newcommand{\NIPPejabatKedua}{}
% Biarkan kosong agar templat hanya menampilkan label dan ruang pengisian.
\newcommand{\TanggalUjian}{}
\newcommand{\TanggalLulus}{}

% --- Data Tesis ---
\newcommand{\JudulTesis}{%
  Judul Karya Ilmiah Maksimum Tiga Baris, Lima Belas Kata Tidak Termasuk
  Kata Depan dan Kata Sambung%
}
\newcommand{\JudulTesisInggris}{%
  Title of Master's Thesis%
}
\newcommand{\KataKunci}{ditulis dalam bahasa Indonesia, disusun berdasarkan abjad, maksimum lima kata atau frasa}
\newcommand{\Keywords}{written in English, arranged alphabetically, maximum five words or phrases}

% Judul ringkasan visual maksimum enam kata, tidak termasuk kata depan/sambung.
\newcommand{\JudulRingkasanVisual}{%
  Judul Ringkasan Visual Maksimum Enam Kata, Tidak Termasuk Kata Depan dan Kata Sambung}
\newcommand{\JudulRingkasanVisualInggris}{%
  Graphical Summary Title Maximum Six Words, Excluding Prepositions and Conjunctions}
% Isi dengan path berkas gambar untuk mengganti kotak placeholder.
\newcommand{\BerkasRingkasanVisual}{}
\newcommand{\BerkasRingkasanVisualInggris}{}

% --- Deklarasi penggunaan AI ---
% Nilai placeholder mempertahankan pilihan [menggunakan/tidak menggunakan]
% seperti pada DOTX. Ganti dengan true atau false setelah menentukan pilihan.
\newcommand{\MenggunakanAI}{placeholder}
\newcommand{\NamaAlatAI}{[nama alat/layanan]}
\newcommand{\TujuanPenggunaanAI}{[alasan]}

% --- Gelar dan sumber dana ---
\newcommand{\GelarMagister}{magister \ldots{}}
% Kosongkan bila penelitian tidak memperoleh pendanaan khusus.
\newcommand{\SumberDana}{}

% --- Data Institusi ---
\newcommand{\Departemen}{Nama Departemen}
\newcommand{\ProgramStudi}{Nama Program Studi}
\newcommand{\ProgramStudiSingkat}{Nama Program Studi}
\newcommand{\Fakultas}{Fakultas/Sekolah}
\newcommand{\Universitas}{Institut Pertanian Bogor}
\newcommand{\KotaUniversitas}{Bogor}
\newcommand{\Tahun}{20XX}

% --- Data Waktu Pelaksanaan (untuk prakata) ---
\newcommand{\BulanMulaiPenelitian}{\ldots{}}
\newcommand{\TahunMulaiPenelitian}{20XX}
\newcommand{\BulanSelesaiPenelitian}{\ldots{}}
\newcommand{\TahunSelesaiPenelitian}{20XX}
\newcommand{\BulanTahunDokumen}{Bulan Tahun}


% Aturan pemenggalan kata
% =============================================================================
% hyphenation-id.tex
%
% Aturan pemenggalan kata (hyphenation) Bahasa Indonesia untuk kata-kata yang
% sering muncul di tesis IPB dan tidak ditangani baik oleh babel default.
%
% Paket: hyphenat (\hyphenation{...}). Tambahkan kata di dalam argumen,
%        pisahkan dengan spasi, gunakan tanda "-" di lokasi pemenggalan.
%
% Catatan: Daftar ini hanya pelengkap; sebagian besar kata sudah ditangani
%          oleh babel[indonesian]. Tambahkan kata yang menurut Anda sering
%          dipenggal salah di sini.
% =============================================================================

\hyphenation{
  agri-kultur
  agro-nomi
  bio-teknolo-gi
  eko-sis-tem
  eko-no-mi
  fito-plank-ton
  herbari-um
  in-sti-tut
  kon-servasi
  ling-kung-an
  manaje-men
  mikro-or-ga-nis-me
  mo-le-kul
  neu-ron
  nu-tri-si
  or-ga-nis-me
  pa-ra-sit
  per-tanian
  per-ubah-an
  pol-u-si
  pop-u-lasi
  sum-ber-daya
  ta-nah
  ta-na-man
  tek-nolo-gi
  varia-bel
  vi-rologi
  zo-oplank-ton
}


% Database referensi untuk biblatex + biber
\addbibresource{src/refs/daftar-pustaka.bib}

% -----------------------------------------------------------------------------
% 2. Page Styles didefinisikan di ipb-thesis.sty
%    - Bagian awal: nomor halaman tidak ditampilkan.
%    - Bagian utama/akhir: nomor Arab di sudut luar atas.
%    - Halaman pembuka bab: nomor tidak ditampilkan.
% -----------------------------------------------------------------------------

% -----------------------------------------------------------------------------
% 3. Metadata untuk hyperref
% -----------------------------------------------------------------------------
\title{\JudulTesis}
\author{\NamaPenulis}
\date{\Tahun}

% -----------------------------------------------------------------------------
% 4. Titik Mulai Dokumen
% -----------------------------------------------------------------------------
\begin{document}

% Hyperref dimuat pada akhir preamble oleh ipb-thesis.sty agar urutan paket
% aman; metadata diisi setelah paket tersebut tersedia.
\hypersetup{
  pdfauthor={\NamaPenulis},
  pdftitle={\JudulTesis},
  pdfsubject={Tesis Magister IPB University -- PPTA 2026},
  pdfkeywords={IPB, tesis, magister, PPTA 2026}
}

\noindent\ignorespaces

% -----------------------------------------------------------------------------
% 5. BAGIAN AWAL
% -----------------------------------------------------------------------------
% Sampul serta halaman pernyataan/hak cipta tidak masuk hitungan nomor
% preliminary pada templat resmi.
\pagenumbering{gobble}
\pagestyle{frontmatter}

% =============================================================================
% Halaman Sampul Digital -- Tesis Magister, PPTA IPB 2026
% =============================================================================

\thispagestyle{empty}
\begin{center}

\vspace*{0.46cm}

\includegraphics[width=2.54cm]{resources/ipb-logo}

\vspace{1.34cm}

{\fontsize{14pt}{17pt}\selectfont\bfseries
\MakeUppercase{\Universitas}\par}

\vspace{1.70cm}

{\fontsize{14pt}{17pt}\selectfont\bfseries
\parbox[t]{12.2cm}{\centering\setlength{\parindent}{0pt}\MakeUppercase{\JudulTesis}}\par}

\vspace{2.01cm}

{\fontsize{14pt}{17pt}\selectfont\bfseries TESIS\par}

\vspace{4.18cm}

{\fontsize{14pt}{17pt}\selectfont\bfseries
\MakeUppercase{\NamaPenulis}\par
\NIM\par}

\vspace{3.65cm}

{\fontsize{14pt}{17pt}\selectfont\bfseries
\MakeUppercase{\Fakultas}\par
\MakeUppercase{\ProgramStudi}\par
\MakeUppercase{\KotaUniversitas}\par
\Tahun\par}

\end{center}

\cleardoublepage
% =============================================================================
% Pernyataan, deklarasi AI, dan pelimpahan hak cipta -- Tesis Magister
% =============================================================================

\ipbprelimheading{PERNYATAAN MENGENAI TESIS,\\
SUMBER INFORMASI, PENGGUNAAN AI, DAN\\
PELIMPAHAN HAK CIPTA}

Dengan ini saya menyatakan bahwa tesis dengan judul
``\JudulTesis'' merupakan salah satu syarat untuk memperoleh gelar
\GelarMagister. Karya ini adalah karya saya dengan arahan dari dosen pembimbing
dan belum diajukan dalam bentuk apa pun kepada perguruan tinggi mana pun.
Sumber informasi yang berasal atau dikutip dari karya yang diterbitkan maupun
tidak diterbitkan dari penulis lain telah disebutkan dalam teks dan dicantumkan
dalam Daftar Pustaka di bagian akhir tesis ini.

\ifdefstring{\MenggunakanAI}{true}{%
  Dalam penyusunan karya ini, saya menggunakan bantuan kecerdasan buatan
  \NamaAlatAI{} untuk \TujuanPenggunaanAI. Setelah menggunakan alat/layanan
  tersebut, saya meninjau dan menyunting konten sesuai kebutuhan serta
  bertanggung jawab penuh atas isi karya tugas akhir ini.
}{%
  \ifdefstring{\MenggunakanAI}{false}{%
    Dalam penyusunan karya ini, saya tidak menggunakan bantuan kecerdasan
    buatan generatif atau teknologi berbantuan AI.
  }{%
    Dalam penyusunan karya ini, saya [menggunakan/tidak menggunakan] bantuan
    kecerdasan buatan \NamaAlatAI{} untuk \TujuanPenggunaanAI. Setelah
    menggunakan alat/layanan tersebut, saya meninjau dan menyunting konten
    sesuai kebutuhan serta bertanggung jawab penuh atas isi karya tugas akhir
    ini.
  }%
}

Dengan ini saya melimpahkan hak cipta dari karya tulis saya kepada Institut
Pertanian Bogor.

\begin{flushright}
\KotaUniversitas, \BulanTahunDokumen\\[1.2cm]
\NamaPenulis\\
\NIM
\end{flushright}

\vfill

\begin{center}
{\fontsize{14pt}{17pt}\selectfont
\copyright{} Hak Cipta milik IPB, \Tahun\\
Hak Cipta dilindungi Undang-Undang\par}
\end{center}

Dilarang mengutip sebagian atau seluruh karya tulis ini tanpa mencantumkan
atau menyebutkan sumbernya. Pengutipan hanya untuk kepentingan pendidikan,
penelitian, penulisan karya ilmiah, penyusunan laporan, penulisan kritik, atau
tinjauan suatu masalah, dan pengutipan tersebut tidak merugikan kepentingan
IPB.

\vspace{10pt}

Dilarang mengumumkan dan memperbanyak sebagian atau seluruh karya tulis ini
dalam bentuk apa pun tanpa izin IPB.

\vspace*{24pt}

\cleardoublepage

% Hitungan Romawi dimulai dari Sorotan (i), tetapi nomornya tidak dicetak.
\frontmatter
\pagestyle{frontmatter}
% =============================================================================
% Sorotan dan Highlights -- Tesis Magister
% - 3--5 poin untuk setiap bahasa.
% - Maksimum 100 karakter per poin, termasuk spasi dan tanda baca.
% - Tidak merujuk pustaka, gambar, atau tabel.
% =============================================================================

\ipbprelimheading{SOROTAN}

\noindent
\MakeUppercase{\NamaPenulis}. \JudulTesis\unskip. Dibimbing oleh
\DaftarPembimbingKapital

\begin{enumerate}[leftmargin=0.5cm,labelwidth=0.25cm,labelsep=0.25cm,
  align=left,itemsep=0pt,parsep=0pt,topsep=12pt,partopsep=0pt]
  \item Penelitian ini membahas [masukkan topik utama], dengan tujuan untuk
    [nyatakan tujuan secara jelas dan ringkas].
  \item Penelitian ini menggunakan [sebutkan metode] untuk
    menganalisis/menginvestigasi [sebutkan objek kajian].
  \item Hasil penelitian menunjukkan bahwa [ringkas temuan utama dengan tepat].
  \item Temuan ini memberikan kontribusi terhadap [bidang/implikasi], serta
    menawarkan wawasan baru mengenai [aspek spesifik].
  \item Tesis ini menunjukkan pentingnya [tekankan signifikansi], dan
    menyarankan potensi penerapan pada [konteks relevan].
\end{enumerate}

\vspace{10pt}

\noindent
Sorotan/\textit{highlights} ditulis dalam satu spasi, disusun 3--5 poin, dan
tidak lebih dari satu halaman. Setiap poin ditulis maksimal 100 karakter
termasuk spasi dan tanda baca. Tidak diperbolehkan mengacu pustaka, gambar,
dan tabel. Singkatan hanya dikenalkan jika masih digunakan lagi dalam bagian
lain sorotan/\textit{highlights}.

\begin{otherlanguage}{english}
\itshape
\vspace{20pt}
\ipbprelimsubheading{HIGHLIGHTS}

\noindent
\MakeUppercase{\NamaPenulisInggris}. \JudulTesisInggris\unskip. Supervised by
\DaftarPembimbingKapitalInggris

\begin{enumerate}[leftmargin=0.5cm,labelwidth=0.25cm,labelsep=0.25cm,
  align=left,itemsep=0pt,parsep=0pt,topsep=7pt,partopsep=0pt]
  \item This study addresses [insert key topic], aiming to [state objective
    clearly and concisely].
  \item The research employs [mention methodology] to analyze/investigate
    [mention subject of study].
  \item Results indicate that [summarize key findings with precision].
  \item The findings contribute to [mention the field/implications], offering
    new insights into [specific aspect].
  \item This thesis demonstrates the importance of [highlight significance],
    suggesting potential applications in [relevant context].
\end{enumerate}
\end{otherlanguage}

\clearpage
% =============================================================================
% Ringkasan Visual dan Graphical Summary -- satu halaman
% Gambar akhir disarankan memiliki resolusi sekurang-kurangnya 300 dpi.
% =============================================================================

\ipbprelimheading{RINGKASAN VISUAL}

\begin{center}
\ifdefempty{\BerkasRingkasanVisual}{%
  \fbox{\parbox[c][5.7cm][c]{0.8\textwidth}{\centering
    Sisipkan sebuah gambar yang secara visual merangkum temuan utama}}%
}{%
  \includegraphics[width=0.8\textwidth,height=5.7cm,keepaspectratio]{\BerkasRingkasanVisual}%
}

\vspace{0.3cm}

{\bfseries\parbox{9cm}{\centering\JudulRingkasanVisual}\par}
\end{center}

\smallskip

\begingroup
\setlength{\parindent}{1.27cm}%
Ringkasan visual/\textit{graphical summary} terdiri atas sebuah gambar yang
secara visual merangkum temuan utama dari tugas akhir. Ringkasan visual ini
disusun secara singkat dan berdiri sendiri, menyampaikan tujuan dan hasil
penelitian, serta disajikan dengan cara yang menarik secara visual dengan
resolusi gambar minimal 300 dpi. Ringkasan visual/\textit{graphical summary}
ditulis dalam satu halaman. Judul diletakkan di bawah ringkasan
visual/\textit{graphical summary}. Mahasiswa wajib mendeklarasikan penggunaan
AI jika digunakan dalam pembuatan ringkasan visual dan perhatikan penggunaan
elemen berhak cipta (lihat Suplemen 6).
\par\endgroup

\begin{otherlanguage}{english}
\itshape
\ipbprelimsubheading{GRAPHICAL SUMMARY}

\begin{center}

\ifdefempty{\BerkasRingkasanVisualInggris}{%
  \fbox{\parbox[c][5.7cm][c]{0.8\textwidth}{\centering
    Insert a single image that visually summarizes the main findings}}%
}{%
  \includegraphics[width=0.8\textwidth,height=5.7cm,keepaspectratio]{\BerkasRingkasanVisualInggris}%
}

\vspace{0.3cm}

{\bfseries\parbox{9cm}{\centering\JudulRingkasanVisualInggris}\par}
\end{center}
\end{otherlanguage}

\clearpage
% =============================================================================
% Ringkasan dan Summary -- Tesis Magister
%
% Berbeda dengan abstrak program sarjana, ringkasan tesis disusun dalam
% beberapa paragraf, dengan spasi tunggal, dan panjangnya tidak lebih dari dua
% halaman. Ringkasan memuat latar belakang permasalahan, tujuan, metode, hasil
% dengan penekanan pada temuan baru, dan implikasi. Tidak diperbolehkan
% mengacu pustaka, gambar, dan tabel.
%
% Ganti teks panduan pada ringkasan-id.tex dan ringkasan-en.tex dengan narasi
% ringkasan yang sebenarnya. Kata kunci diatur di src/config/information.tex.
% =============================================================================

\ipbprelimheading{RINGKASAN}
% Ganti teks panduan ini dengan narasi ringkasan dalam bahasa Indonesia.
% Ringkasan disusun dalam beberapa paragraf, spasi tunggal, maksimum dua
% halaman, dan tidak mengacu pustaka, gambar, atau tabel.
\noindent
\MakeUppercase{\NamaPenulis}. \JudulTesis\unskip. Dibimbing oleh
\DaftarPembimbingKapital

Narasi ditulis dalam satu spasi, disusun dalam beberapa paragraf, dan tidak
lebih dari dua halaman. Ringkasan/\textit{summary} memuat latar belakang
permasalahan, tujuan, metode, hasil dengan penekanan pada temuan baru, dan
implikasi yang disajikan secara ringkas, informatif, dan faktual. Tidak
diperbolehkan mengacu pustaka, gambar, dan tabel. Singkatan hanya dikenalkan
jika masih digunakan lagi dalam bagian lain dari ringkasan/\textit{summary}.
Ringkasan dalam bahasa Inggris ditulis dengan huruf miring
(\textit{italic}).


\ipbprelimsubheading{\itshape SUMMARY}
% Replace this guidance text with the English summary narrative.
% The summary is written in several paragraphs, single-spaced, at most two
% pages, and must not refer to literature, figures, or tables.
\begin{otherlanguage}{english}
\itshape
\noindent
\MakeUppercase{\NamaPenulisInggris}. \JudulTesisInggris\unskip. Supervised by
\DaftarPembimbingKapitalInggris

The narrative is written in single spacing, arranged in several paragraphs,
and not more than two pages. The summary contains the background of the
problem, objectives, methods, results with emphasis on new findings, and
implications presented concisely, informatively, and factually. It must not
refer to literature, figures, or tables. Abbreviations are introduced only if
they are used again in another part of the summary. The English summary is
written in italics.
\end{otherlanguage}


\vspace{0.48cm}

\noindent\hangindent=2cm\hangafter=1
Kata kunci: \KataKunci

\begin{otherlanguage}{english}
\itshape
\noindent\hangindent=2cm\hangafter=1
Keywords: \Keywords
\par
\end{otherlanguage}

\clearpage
% =============================================================================
% Tim Penguji pada Ujian Tesis
% =============================================================================

\thispagestyle{empty}
\vspace*{\fill}

\noindent Tim Penguji pada Ujian Tesis:

\begin{tabular}{@{}r@{\hspace{0.7cm}}p{0.8\textwidth}@{}}
1 & \NamaPengujiSatu \\
2 & \NamaPengujiDua
\end{tabular}

\vspace{0.3cm}

\clearpage
% =============================================================================
% Halaman Pengesahan -- Tesis Magister, PPTA IPB 2026
%
% Memuat daftar tim penguji, data tesis, tanda tangan komisi pembimbing,
% Ketua Program Studi, dan Dekan Fakultas/Sekolah.
% Pembimbing ketiga hanya ditampilkan bila \TampilkanPembimbingTiga = true
% (lihat src/config/information.tex).
% =============================================================================

\thispagestyle{empty}

\noindent
\begin{tabular}{@{}p{2.5cm}@{}p{0.2cm}@{\hspace{0.03cm}}p{0.685\textwidth}@{}}
Judul Tesis & : & \JudulTesis \\
Nama        & : & \NamaPenulis \\
NIM         & : & \NIM
\end{tabular}

\vspace{2.447cm}

\begin{center}
\hspace*{1cm}Disetujui oleh
\end{center}

\vspace{1.146cm}

% Garis tanda tangan diletakkan di antara baris jabatan dan baris nama, sesuai
% templat Word (garis berada pada baris pertama, nama pada baris berikutnya).
\noindent
\begin{tabular}{p{8.33cm}>{\raggedleft\arraybackslash}p{4.08cm}}
Pembimbing 1:\newline
\NamaPembimbingSatu & \raisebox{-4.7pt}[0pt][0pt]{\rule{4cm}{0.4pt}}\newline\mbox{} \\[1.95cm]
Pembimbing 2:\newline
\NamaPembimbingDua & \raisebox{-4.7pt}[0pt][0pt]{\rule{4cm}{0.4pt}}\newline\mbox{}%
\ifdefstring{\TampilkanPembimbingTiga}{true}{\\[1.95cm]
Pembimbing 3:\newline
\NamaPembimbingTiga & \raisebox{-4.7pt}[0pt][0pt]{\rule{4cm}{0.4pt}}\newline\mbox{}}{}
\end{tabular}

\vspace{0.926cm}

\begin{center}
\hspace*{1cm}Diketahui oleh
\end{center}

\vspace{1.035cm}

\noindent
\begin{tabular}{p{8.08cm}>{\raggedleft\arraybackslash}p{4.08cm}}
Ketua Program Studi:\newline
\NamaKetuaProgramStudi\newline
NIP \NIPKetuaProgramStudi & \mbox{}\newline\rule{4cm}{0.4pt}\newline\mbox{} \\[2.03cm]
\JabatanPejabatKedua:\newline
\NamaPejabatKedua\newline
NIP \NIPPejabatKedua & \mbox{}\newline\rule{4cm}{0.4pt}\newline\mbox{}
\end{tabular}

\vfill

\noindent
\begin{tabular*}{\textwidth}{@{\extracolsep{\fill}}ll@{}}
Tanggal Ujian: \TanggalUjian & Tanggal Lulus: \TanggalLulus
\end{tabular*}
\par
\vspace*{9pt}

\cleardoublepage
% =============================================================================
% Prakata (Kata Pengantar) -- Tesis Magister
%
% PPTA 2026: Persantunan diberikan kepada pihak yang berkontribusi langsung,
% termasuk sumber dana bila ada.
% =============================================================================

\ipbprelimheading{PRAKATA}

Puji dan syukur penulis panjatkan kepada Tuhan Yang Maha Esa atas segala
karunia-Nya sehingga tesis ini berhasil diselesaikan. Judul dalam penelitian
yang dilaksanakan sejak bulan \BulanMulaiPenelitian{} \TahunMulaiPenelitian{}
sampai bulan \BulanSelesaiPenelitian{} \TahunSelesaiPenelitian{} ini ialah
``\JudulTesis''.

Terima kasih penulis ucapkan kepada para pembimbing,
\DaftarPembimbing, yang telah membimbing
dan banyak memberi saran. Ucapan terima kasih juga disampaikan kepada
pembimbing akademik, moderator seminar, dan penguji luar komisi pembimbing.
Penghargaan penulis sampaikan kepada \ldots{} (nama lengkap dan gelar dari
lembaga/instansi/perusahaan yang telah memberi izin penelitian), \ldots{}
(nama dan gelar atau bapak/ibu jika tidak ada gelar) beserta staf Laboratorium
\ldots{} yang telah membantu selama pengumpulan data. Ungkapan terima kasih
juga disampaikan kepada ayah, ibu, serta seluruh keluarga yang telah
memberikan dukungan, doa, dan kasih sayangnya \ldots{} dan seterusnya.

\ifdefempty{\SumberDana}{}{Penelitian ini didukung oleh \SumberDana.}

Semoga tugas akhir ini bermanfaat bagi pihak yang membutuhkan dan bagi
kemajuan ilmu pengetahuan.

\vspace{14pt}

\noindent\textcolor{red}{Catatan: Ucapan terima kasih diberikan hanya kepada
pihak-pihak yang secara langsung berkontribusi terhadap pengumpulan data dan
penulisan tugas akhir.}

\vspace{2\baselineskip}

\begin{flushright}
\KotaUniversitas, \BulanTahunDokumen\\[\baselineskip]
\NamaPenulis
\end{flushright}

\cleardoublepage
\input{preliminaries/daftar-isi}

% -----------------------------------------------------------------------------
% 6. BAGIAN INTI (BAB) -- Bagian Inti / Pokok
%    - Nomor halaman: Arab (1, 2, 3, ...), sudut luar atas
%    - Margin fisik: sisi dalam 4 cm, sisi luar 3 cm, atas/bawah 3 cm
% -----------------------------------------------------------------------------
\mainmatter
\pagestyle{mainmatter}

% =============================================================================
% BAB I PENDAHULUAN -- teks panduan templat format tesis terbaru
% Ganti seluruh teks panduan dengan isi tesis.
% =============================================================================

\chapter{PENDAHULUAN}

Bab pendahuluan memuat latar belakang atau justifikasi dipilihnya topik karya
ilmiah tugas akhir, perumusan atau pendekatan penyelesaian masalah, tujuan,
manfaat, dan ruang lingkup. Di dalam pendahuluan dijelaskan perumusan atau
pendekatan penyelesaian masalah dan alasan pemilihan metode yang digunakan.
Merujuk pada proses perumusan masalah tugas akhir, bagian Kerangka Pikir dan
Hipotesis dapat ditulis di sini, tidak ditulis dalam bab tersendiri. Untuk
membantu mengikuti alur pikir secara skematis, dapat juga dibuat bagan alir
kerangka proses dan rumusan masalah serta pencapaian tujuan tugas akhir.

\section{Latar Belakang}

Latar belakang berisi penjelasan alasan memilih topik dan pentingnya kajian
tugas akhir itu dilakukan berdasarkan alasan teoretis dan praktis, serta
bagaimana masalah tersebut dapat diatasi dan manfaat dari penyelesaian masalah.
Paparan tidak berbelit-belit atau dimulai dengan latar belakang yang umum.
Pernyataan mengenai apa yang diteliti dan apa yang diharapkan diawali dengan
pemikiran logis. Pemaparan latar belakang harus sistematis, logis, serta
disertai data, informasi, dan telaah pustaka dari sumber primer, mutakhir, dan
relevan yang dapat dipertanggungjawabkan secara ilmiah. Masalah penelitian
yang lebih spesifik dirumuskan pada bagian rumusan masalah.

\section{Rumusan Masalah}

Rumusan masalah merupakan pernyataan ringkas mengenai masalah yang akan
diselesaikan dan cara mengatasinya untuk menjawab tujuan tugas akhir. Masalah
yang diteliti dapat dirumuskan karena berbagai sebab, seperti adanya
kesenjangan (\textit{gap}), tantangan, kesangsian, ketidakjelasan, dan
keingintahuan secara akademik yang berkaitan dengan fenomena alam, sosial, dan
ekonomi. Dalam pernyataan ringkas tersebut harus tercakup pendekatan yang
digunakan dalam perumusan masalah.

\section{Tujuan}

Pernyataan tujuan tugas akhir ialah pernyataan singkat dan jelas tentang tujuan
yang akan dicapai sebagai upaya pemecahan masalah maupun memahami gejala
(fenomena) yang dijelaskan dalam latar belakang. Tujuan merupakan pemandu atau
arah untuk merencanakan dan melaksanakan kajian tugas akhir. Tujuan penelitian
ditulis dengan memilih kata kerja yang hasilnya dapat diukur dan dilihat,
seperti: menguraikan, menerangkan, membuktikan, menjajaki, menguji,
membuktikan, atau menerapkan suatu gejala, konsep atau dugaan, atau bahkan
membuat suatu prototipe. Jangan menggunakan kata kerja mengetahui atau
memahami. Masalah dan tujuan penelitian harus terkait dan konsisten.

\section{Manfaat}

Manfaat merupakan dampak positif (kegunaan) dari hasil karya ilmiah tugas akhir
bagi bidang ipteks, pembangunan, dan masyarakat. Manfaat utama dari hasil tugas
akhir adalah menambah khasanah ilmu pengetahuan dalam bentuk pustaka sebagai
sumber acuan/referensi untuk pengembangan ipteks, para pengambil keputusan baik
di industri maupun pemerintah dan lembaga untuk menyusun kebijakan baru, serta
masyarakat umum. Manfaat dinyatakan dengan kata kerja yang lugas dan logis.

\section{Ruang Lingkup (opsional)}

Tugas akhir sering kali dihadapkan pada keterbatasan data, dana, waktu, metode,
bahkan teori. Oleh karena itu, tugas akhir perlu dengan tegas menunjukkan ruang
lingkup dengan mempertimbangkan keterbatasan tersebut.

\section{Hipotesis (opsional)}

Hipotesis dapat ditulis secara eksplisit atau tersirat sesuai bidang ipteks yang
relevan. Hipotesis dapat menjadi bagian dari pendahuluan (untuk bidang sains,
teknik, dan kesehatan), atau merupakan bagian akhir dari tinjauan pustaka
(untuk bidang sosial, ekonomi, dan humaniora).
    % BAB I   PENDAHULUAN
% =============================================================================
% BAB II TINJAUAN PUSTAKA -- teks panduan templat format tesis terbaru
% =============================================================================

\chapter{TINJAUAN PUSTAKA}

Tinjauan pustaka berisi telaah/ulasan atas pustaka-pustaka yang relevan dengan
topik tugas akhir untuk mendapatkan informasi yang lengkap terkait kemajuan
ipteks yang telah diketahui sampai yang terkini (\textit{state of the art}).
Hal ini untuk meyakinkan pembaca bahwa tugas akhir yang dilaporkan adalah
pengetahuan baru yang lebih maju dari pengetahuan sebelumnya. Pustaka yang
digunakan dalam bab ini ialah acuan primer, diutamakan artikel jurnal dan paten
yang relevan dengan bidang yang diteliti, terkini, dan asli. Pustaka acuan harus
kredibel, dan mutakhir (setidaknya 80\% dalam 10 tahun terakhir). Diktat dan
buku ajar tidak termasuk acuan primer. Referensi lama berupa teori umum/mapan
tetap dapat digunakan. Tinjauan pustaka ditulis dengan ketentuan, jumlah halaman
bab tersebut tidak melebihi 10\% dari total halaman bagian utama naskah dan
tidak melebihi jumlah halaman Hasil dan Pembahasan. Tinjauan pustaka memuat
telaah singkat, jelas, dan sistematis tentang kerangka teoretis, kerangka pikir,
temuan, postulat-postulat, prinsip, asumsi, dan hasil penelitian yang relevan
yang melandasi topik tugas akhir atau gagasan guna menggali pemahaman mengenai
masalah kajian tugas akhir dan pemecahan masalahnya. Dalam penelitian bidang
ilmu sosial dan ekonomi, tinjauan pustaka menjadi dasar penyusunan kerangka
analisis baru dan hipotesis baru dalam topik karya ilmiah tersebut.

\section[Judul Subbab]{Judul Subbab (Kata dalam judul diawali huruf kapital dan dicetak tebal)}

Uraian dengan deskripsi untuk judul subbab 1 \dotfill

\noindent\dotfill

\noindent\dotfill

\begin{enumerate}[label=\alph*,leftmargin=0.75cm,labelwidth=0.5cm,
  labelsep=0.25cm,align=left,itemsep=0pt,parsep=0pt,topsep=0pt,partopsep=0pt]
  \item \ldots
  \item \ldots
  \item \ldots
  \begin{enumerate}[label=\arabic*),leftmargin=0.75cm,labelwidth=0.5cm,
    labelsep=0.25cm,align=left,itemsep=0pt,parsep=0pt,topsep=0pt,partopsep=0pt]
    \item \ldots
    \item \ldots
    \item \ldots
  \end{enumerate}
\end{enumerate}

\subsection{Judul Sub-subbab (Kata dalam judul diawali huruf kapital dan dicetak tidak tebal)}

\begin{ipbsubsubcontent}
Berikut adalah contoh uraian dengan deskripsi pada sub-subbab. Pada sub-subbab
ini posisi paragraf lebih menjorok 0,5 cm dari paragraf di subbab. \dotfill

\noindent\dotfill

\noindent\dotfill

\noindent\dotfill

\begin{enumerate}[label=\alph*),leftmargin=0.75cm,labelwidth=0.5cm,
  labelsep=0.25cm,align=left,itemsep=0pt,parsep=0pt,topsep=0pt,partopsep=0pt]
  \item \ldots
  \item \ldots
  \item \ldots
  \begin{enumerate}[label=(\arabic*),leftmargin=0.75cm,labelwidth=0.5cm,
    labelsep=0.25cm,align=left,itemsep=0pt,parsep=0pt,topsep=0pt,partopsep=0pt]
    \item \ldots
    \item \ldots
    \item \ldots
  \end{enumerate}
\end{enumerate}
\end{ipbsubsubcontent}

\subsection{Judul Sub-subbab (Kata dalam judul diawali huruf kapital dan dicetak tidak tebal)}

\begin{ipbsubsubcontent}
Berikut adalah contoh uraian dengan deskripsi pada sub-subbab. Pada sub-subbab
ini posisi paragraf lebih menjorok 0,5 cm dari paragraf di subbab. \dotfill

\noindent\dotfill
\end{ipbsubsubcontent}

\section[Judul Subbab]{Judul Subbab (Kata dalam judul diawali huruf kapital dan dicetak tebal)}

Berikut adalah contoh uraian dengan deskripsi pada subbab. \dotfill

\noindent\dotfill

\noindent\dotfill

\section[Judul Subbab]{Judul Subbab (Kata dalam judul diawali huruf kapital dan dicetak tebal)}

Berikut adalah contoh uraian dengan deskripsi pada subbab. \dotfill

\noindent\dotfill

\noindent\dotfill

\noindent\dotfill

\noindent\dotfill\ (Gambar~\ref{fig:mukosa-trenggiling}).

\begin{figure}[htbp]
  \centering
  \vspace{0.25cm}
  \fbox{\parbox[c][9.8cm][c]{0.92\textwidth}{\centering
    Sisipkan mikrograf atau gambar hasil penelitian di sini}}
  \caption[Mikrograf mukosa lambung trenggiling
    (\textit{Manis javanica})]{Mikrograf mukosa lambung trenggiling
    (\textit{Manis javanica}). (A) Seluruh permukaan mukosa lambung
    trenggiling dilapisi oleh epitel pipih banyak-lapis yang mengalami
    keratinisasi. (B) Pada lubang kelenjar, epitel berubah menjadi epitel
    silindris (tanda panah). Lapisan epitel mukosa terdiri atas: korneum (a),
    granulosum (tanda panah), spinosum (b), dan germinativum (c). Lapis
    submucosa (d) terletak di profundal lapis germinativum. Pewarnaan HE,
    skala = 50~$\upmu$m. (dimodifikasi dari Nisa' \textit{et al.} (2010),
    \textit{Anatomia Histologia Embryologia}. 39:432--439).}
  \label{fig:mukosa-trenggiling}
\end{figure}
    % BAB II  TINJAUAN PUSTAKA
% =============================================================================
% BAB III METODE -- teks panduan templat format tesis terbaru
% =============================================================================

\chapter{METODE}

Bab ini dapat diawali dengan kerangka pendekatan studi. Metode dapat berupa
percobaan laboratorium, percobaan lapangan, dan survei lapangan yang dirancang
sesuai dengan tujuan atau jenis penelitian, seperti: eksploratif, deskriptif,
koreksional, kausal, komparatif, eksperimen, tindakan
(\textit{action research}), pemodelan, analisis suatu teori, atau kombinasi
dari berbagai jenis penelitian. Judul Subbab pada Metode disesuaikan dengan
bidang keilmuan. Untuk tugas akhir yang menggunakan metode kualitatif, jelaskan
pendekatan yang digunakan, proses pengumpulan dan analisis informasi, dan
proses penafsiran hasil tugas akhir. Maksud dari perincian ini ialah untuk
menjamin keterulangan hasil. Berikut contoh subbab metode penelitian.

\section[Waktu dan Tempat]{Waktu dan Tempat (Judul subbab disesuaikan dengan bidang keilmuan)}

Deskripsi\dotfill

\noindent\dotfill

\noindent\dotfill

\section[Alat dan Bahan]{Alat dan Bahan (Judul subbab disesuaikan dengan bidang keilmuan)}

Deskripsi\dotfill

\noindent\dotfill

\noindent\dotfill

\section[Prosedur Kerja]{Prosedur Kerja (Judul subbab disesuaikan dengan bidang keilmuan)}

\subsection{Judul Sub-subbab (Kata dalam judul diawali huruf kapital dan dicetak tidak tebal)}

\begin{ipbsubsubcontent}
Berikut adalah contoh uraian dengan deskripsi pada sub-subbab. Pada sub-subbab
ini posisi paragraf lebih menjorok 0,5 cm dari paragraf di subbab. \dotfill

\noindent\dotfill

\noindent\dotfill

\noindent\dotfill

\noindent\dotfill

\noindent\dotfill

\noindent Contoh penulisan formula

\begin{equation*}
  f(x)=a_{0}+\sum_{n=1}^{\infty}
  \left(a_{n}\cos\frac{n\pi x}{L}
  +b_{n}\sin\frac{n\pi x}{L}\right)
\end{equation*}

\noindent Keterangan:

\noindent Berikan deskripsi pada setiap notasi formula di atas.
\end{ipbsubsubcontent}

\section{Analisis Data}

Deskripsi\dotfill

\noindent\dotfill

\noindent\dotfill
    % BAB III METODE
% =============================================================================
% BAB IV HASIL DAN PEMBAHASAN -- teks dan contoh dari templat format tesis terbaru
% =============================================================================

\chapter{HASIL DAN PEMBAHASAN (terpisah atau gabung)}

Dalam penulisan hasil dan pembahasan dapat dipisah sebagai bab Hasil dan bab
Pembahasan, atau digabung menjadi bab Hasil dan Pembahasan. Bila Hasil dan
Pembahasan disatukan dalam satu bab, data/temuan disajikan dahulu, diberi
penjelasan yang cukup untuk temuan yang penting, dilanjutkan dengan penafsiran,
kemudian dengan pembahasan.

\section{Hasil}

Hasil menampilkan data/temuan dan dapat dibagi dalam beberapa subbab sesuai
dengan tujuan. Data dapat disajikan dengan ilustrasi dalam bentuk Tabel atau
Gambar (peta, denah, foto, diagram). Ilustrasi harus mampu berdiri sendiri,
artinya mudah dipahami pembaca tanpa harus merujuk teks. Tujuannya adalah
membantu pembaca memahami data dan menarik informasi penting dengan cepat
(\textit{storytelling with data}). Semua ilustrasi harus diletakkan
sedekat-dekatnya dengan teks yang menyatakan keberadaannya. Perujukan pada
ilustrasi dinyatakan di dalam paragraf sebelum tabel atau gambar. Kata
\enquote{tabel} dan \enquote{gambar} diawali dengan huruf kapital bila diikuti
dengan nomor. Nomor tabel atau gambar berurut sesuai dengan urutan kemunculannya
dalam tubuh tulisan. Contohnya adalah sebagai berikut:

\noindent\ldots{} seperti ditunjukkan pada Gambar 5.\par
\noindent\ldots{} mendekati bentuk sigmoid (Gambar 5).\par
\noindent\ldots{} meningkat dengan pesat (Tabel 3).\par
\noindent\ldots{} (lihat Lampiran 1).\par

Perujukan pada ilustrasi yang tidak disertai dengan keterangan harus
dihindari.\par
\noindent Misalnya:\par
\noindent \enquote{Hasilnya dapat dilihat pada Tabel 1} atau
\enquote{Hasilnya disajikan pada Tabel 1}.\par
\noindent Pernyataan yang lebih baik ialah\par
\noindent \enquote{Pada Tabel 1 ditunjukkan bahwa kacang hijau lebih banyak
memancarkan spektrum biru daripada kacang tanah}.

\vspace{-4pt}
\begin{table}[htbp]
  \captionsetup{justification=centering,singlelinecheck=true}
  \caption[Judul tabel biasanya pendek tanpa diakhiri tanda titik]%
    {Judul tabel biasanya pendek tanpa diakhiri tanda titik$^{\mathrm{a}}$}
  \label{tab:contoh-struktur}
  {\fontsize{11pt}{13pt}\selectfont
  \renewcommand{\arraystretch}{1.15}%
  \begin{tabular}{@{}p{3.47cm}@{}*{4}{>{\centering\arraybackslash}p{2.63cm}@{}}}
    \toprule
    \multirow{2}{3.47cm}{\centering\textbf{Judul kolom\newline pertama}}
      & \multicolumn{2}{c}{\textbf{Judul kolom$^{\mathrm{b}}$}}
      & \multicolumn{2}{c}{\textbf{Judul kolom$^{\mathrm{b}}$}} \\
    \cmidrule(lr){2-3}\cmidrule(l){4-5}
      & \hspace{5.2pt}\textbf{Subjudul\newline kolom}
      & \hspace{5.2pt}\textbf{Subjudul\newline kolom}
      & \hspace{5.2pt}\textbf{Subjudul\newline kolom}
      & \hspace{5.2pt}\textbf{Subjudul\newline kolom} \\
    \specialrule{\lightrulewidth}{0pt}{2pt}
    & \multicolumn{4}{c}{[area informasi]} \\
    Judul baris$^{\mathrm{c}}$ & & & & \\
    \quad Subjudul baris & & & & \\
    \quad Subjudul baris & & & & \\
    Judul baris & & & & \\
    \quad Subjudul baris & & & & \\
    \quad Subjudul baris & & & & \\
    & \multicolumn{4}{c}{[area informasi]} \\
    \bottomrule
  \end{tabular}
  }
  \ipbtablenote{$^{\mathrm{a}}$[catatan kaki] [titik] Sumber [jika ada]: xxxx
    xxxx [titik] $^{\mathrm{b}}$[catatan kaki] xxxx xxxx [titik]
    $^{\mathrm{c}}$[catatan kaki] xxxx xxxx [titik]}
\end{table}

\noindent Catatan kaki diketik dengan font Times New Roman 10 pt. Antara baris pertama
catatan kaki dan garis dasar tabel diberi jarak 3 pt.

\clearpage

\begin{table}[H]
  \captionsetup{format=ipbhangonehalf,justification=justified,
    singlelinecheck=false}
  \caption[Tingkat kekerasan dan kandungan gula buah pisang ambon pada suhu
    simpan yang berbeda dan pemberian putresina]{Tingkat kekerasan dan
    kandungan gula buah pisang ambon pada suhu
    simpan yang berbeda dan pemberian putresina
    \textbf{(Jika judul tabel lebih dari satu baris, maka baris selanjutnya
    sejajar dengan kata pertama pada judul tabel)}}
  \label{tab:pisang-ambon}
  {\fontsize{11pt}{12pt}\selectfont
  \begin{tabularx}{13.63cm}{@{}Xccc@{}}
    \toprule
    \multirow{2}{*}{\textbf{Perlakuan}}
      & \multicolumn{3}{c}{\textbf{Kekerasan buah dan kandungan gula pada hari ke-}} \\
    \cmidrule(l){2-4}
      & \textbf{0} & \textbf{7} & \textbf{14} \\
    \midrule
    & \multicolumn{3}{c}{Kekerasan buah (mm 50 g$^{-1}$ detik$^{-1}$)$^{\mathrm{a}}$} \\
    Suhu simpan & & & \\
    \quad $15\,^{\circ}\mathrm{C}$ & 10,20$^{\mathrm{a}}$ & 13,40$^{\mathrm{a}}$ & 11,83$^{\mathrm{a}}$ \\
    \quad $28\,^{\circ}\mathrm{C}$ & 10,64$^{\mathrm{a}}$ & 14,22$^{\mathrm{a}}$ & 88,43$^{\mathrm{b}}$ \\
    Putresina & & & \\
    \quad Dengan putresina & 11,07$^{\mathrm{a}}$ & 13,23$^{\mathrm{a}}$ & 21,19$^{\mathrm{a}}$ \\
    \quad Tanpa putresina & 10,76$^{\mathrm{a}}$ & 14,40$^{\mathrm{a}}$ & 41,82$^{\mathrm{b}}$ \\
    \midrule
    & \multicolumn{3}{c}{Gula (\%)$^{\mathrm{a}}$} \\
    Suhu simpan & & & \\
    \quad $15\,^{\circ}\mathrm{C}$ & 0,38$^{\mathrm{a}}$ & 0,56$^{\mathrm{a}}$ & 0,73$^{\mathrm{a}}$ \\
    \quad $28\,^{\circ}\mathrm{C}$ & 0,55$^{\mathrm{a}}$ & 1,82$^{\mathrm{a}}$ & 14,41$^{\mathrm{b}}$ \\
    Putresina & & & \\
    \quad Dengan putresina & 0,53$^{\mathrm{a}}$ & 0,87$^{\mathrm{a}}$ & 6,98$^{\mathrm{a}}$ \\
    \quad Tanpa putresina & 0,40$^{\mathrm{a}}$ & 1,52$^{\mathrm{a}}$ & 6,91$^{\mathrm{a}}$ \\
    \bottomrule
  \end{tabularx}
  }
  \ipbtablenote{$^{\mathrm{a}}$Angka-angka pada kolom yang sama yang diikuti
    oleh huruf yang sama tidak berbeda nyata pada taraf uji 5\% (uji selang
    berganda Duncan).}
\end{table}

\begin{figure}[H]
  \centering
  \vspace{0.65cm}
  \fbox{\parbox[c][6.4cm][c]{0.92\textwidth}{\centering
    Sisipkan diagram batang; gunakan warna dan motif yang tetap terbaca}}
  \caption[Pemilihan warna dan motif batang pada diagram batang yang jelas
    menunjukkan keterbacaan \textit{error bar}]{Pemilihan warna dan motif
    batang pada diagram batang yang jelas
    menunjukkan keterbacaan \textit{error bar}, (a) dua batang dengan penanda
    signifikansi (*), (b) beberapa batang dengan variasi warna yang ramah buta
    warna dan motif yang berbeda namun mudah dibaca. Penanda statistika
    (asteriks atau huruf) dapat diletakkan langsung di atas balok. Notasi pada
    tiap batang sudah jelas sehingga legenda tidak dicantumkan. Legenda atau
    keterangan lain dapat dicantumkan di judul gambar. Data pada diagram adalah
    \textit{random} dan dibuat dengan Microsoft Excel for Mac Version 16.89.}
  \label{fig:diagram-batang}
\end{figure}

\clearpage

Penempatan ilustrasi (tabel dan gambar) di dalam teks memerlukan perhatian
khusus untuk memastikan kualitas visualisasi ilmiah, keterbacaan, dan estetika.
Pengaturan yang tepat menjamin pembaca dapat dengan mudah menghubungkan teks
dengan ilustrasi yang disajikan. Berikut adalah panduan umum tentang penempatan
ilustrasi.

\begin{enumerate}[label=\alph*.,leftmargin=0.75cm,labelwidth=0.5cm,
  labelsep=0.25cm,align=left,itemsep=0pt,parsep=0pt,topsep=0pt,partopsep=0pt]
  \item Kedekatan dengan teks rujukan: Setiap ilustrasi ditempatkan
    sedekat mungkin setelah kalimat di dalam teks yang merujuk atau membahas
    ilustrasi tersebut. Hal ini akan meminimalkan gangguan alur baca dan
    memudahkan pembaca mengaitkan informasi.
  \item Posisi peletakan ideal: Ilustrasi diletakkan di tengah
    halaman secara horizontal, memanfaatkan ruang yang tersedia. Ilustrasi
    diposisikan sedemikian rupa sehingga judul ilustrasi berada di satu halaman
    yang sama. Penting juga untuk tidak meninggalkan ruang kosong setelah
    ilustrasi dan memastikan ada teks/paragraf yang mengikutinya.
  \item Efisiensi ruang dan orientasi: Ilustrasi disajikan dan
    diatur agar muat dalam satu halaman. Jika diperlukan, orientasi halaman
    bisa diubah menjadi lanskap untuk mengakomodasi ilustrasi yang lebar.
    Untuk tabel, format yang pendek dan lebar lebih disukai daripada yang
    panjang dan sempit. Apabila tabel dengan baris yang banyak tidak dapat
    dihindari, judul kolom berulang (\textit{repeat header row}) bisa
    digunakan, sehingga tiap halaman sambungan memiliki judul kolom.
\end{enumerate}

\vspace{13pt}

\section{Pembahasan}

Pembahasan merupakan interpretasi atau penjelasan atas data hasil kegiatan
tugas akhir. Dalam pembahasan, atau penjelasan tersebut harus dikaitkan dengan
pustaka, terutama yang mutakhir, primer, dan kredibel. Pernyataan-pernyataan
dalam paragraf dikemas dengan baik, dimulai dari pendapat sendiri di awal
paragraf, diikuti dengan dukungan pustaka, dan diakhiri dengan kalimat
penyimpulan. Argumentasi dikemukakan secara singkat dan logis yang difokuskan
untuk menjawab tujuan dan menguji hipotesis (jika ada). Argumentasi dikemukakan
untuk menunjukkan persamaan, membahas perbedaan, dan penyebab timbulnya
perbedaan, dikaitkan dengan teori dan temuan-temuan sebelumnya yang relevan.
Pembahasan berupaya menunjukkan aspek-aspek baru yang ditemukan dan merupakan
satu kesatuan. Pembahasan juga dapat mengemukakan keterbatasan dalam kegiatan
tugas akhir yang dilaksanakan.

Terdapat tiga tingkatan kedalaman dalam pembahasan hasil penelitian, yaitu
menunjukkan kejadian, pola, dan struktur. Dalam menunjukkan kejadian, peneliti
hanya membaca fakta-fakta hasil penelitian/kajian yang diperoleh. Dalam
menunjukkan pola, peneliti memperhatikan keteraturan atau kecenderungan dari
fakta-fakta hasil penelitian/kajian. Dalam menunjukkan struktur, peneliti
menjelaskan lebih lanjut dari pola dengan menjawab pertanyaan mengapa pola itu
terjadi, dan faktor apa yang menyebabkan pola itu terjadi. Pembahasan dengan
menunjukkan struktur memerlukan dukungan teori atau rujukan hasil-hasil
penelitian/kajian terdahulu yang telah teruji secara ilmiah. Pendapat peneliti
terdahulu yang sudah diringkas dalam bab Pendahuluan atau bab Tinjauan Pustaka
tidak perlu diulang lagi, tetapi diacu saja seperlunya. Sebaiknya, dalam
pembahasan juga dapat diberikan implikasi penerapan temuan baru dan segi-segi
lain yang perlu diteliti lebih lanjut. Pembahasan diakhiri dengan kalimat
positif, tegas, dan kuat.
    % BAB IV  HASIL DAN PEMBAHASAN
% =============================================================================
% BAB V SIMPULAN DAN SARAN -- teks panduan templat format tesis terbaru
% =============================================================================

\chapter{SIMPULAN DAN SARAN}

\section{Simpulan}

Simpulan ditulis dalam bentuk paragraf yang efektif sesuai dengan tujuan
penelitian. Simpulan merupakan jawaban dari tujuan yang sudah ditentukan dan
tidak dimaksudkan sebagai ringkasan hasil. Simpulan merupakan hasil penelitian
yang boleh jadi telah dikemukakan dalam perumusan masalah dan telah diberi
jawaban sementara berupa hipotesis. Dalam menulis simpulan, penulis harus
membedakan dugaan, temuan, dan simpulan hasil studi. Pernyataan simpulan harus
dilakukan secara cermat dan hati-hati. Penyampaian simpulan ini dapat dilakukan
sebanyak 3 kali, yakni dalam pembahasan, simpulan, dan abstrak sehingga
diperlukan kecermatan untuk menyajikannya dengan ungkapan yang berbeda-beda.

\section{Saran}

Saran sebaiknya mengarah ke implikasi atau tindakan lanjutan yang harus
dilakukan sehubungan dengan temuan atau simpulan penulis. Saran yang
dikemukakan harus berkaitan dengan pelaksanaan atau hasil penelitian. Dengan
demikian saran ini mengemukakan hal-hal yang perlu diteliti lebih lanjut
terutama untuk memperbaiki kelemahan atau kekurangan dalam penelitian yang
dilakukan atau perbaikan asumsi yang diambil sehingga didapatkan hasil yang
lebih baik. Jadi, saran tersebut harus diuraikan secara spesifik. Jangan
menyarankan hal-hal yang tidak dianalisis dan dibahas dalam penelitian serta
terkesan menggurui atau memuaskan keinginan peneliti. Untuk penelitian yang
berkaitan dengan permasalahan kebijakan, tidak perlu menyarankan kebijakan yang
tidak berkaitan dengan hasil penelitian.
    % BAB V   SIMPULAN DAN SARAN

% -----------------------------------------------------------------------------
% 7. BAGIAN AKHIR
%    - Nomor halaman: Arab (lanjut), sudut luar atas
% -----------------------------------------------------------------------------
\backmatter
\pagestyle{mainmatter}

\input{backmatter/daftar-pustaka}
% =============================================================================
% Lampiran
%
% PPTA 2026: Diawali halaman pemisah. Setiap lampiran di halaman baru.
% Caption lampiran menggunakan
% format "Lampiran N  Judul" dengan hanging indent. Counter lampiran
% terpisah dari counter chapter dan figure/table.
% =============================================================================

\cleardoublepage
\thispagestyle{empty}
% Posisi judul mengikuti templat Word: sedikit di atas titik tengah halaman.
\vspace*{262.4pt}
\begin{center}
{\fontsize{14pt}{17pt}\selectfont\bfseries LAMPIRAN}
\end{center}
\addcontentsline{toc}{chapter}{LAMPIRAN}
\vspace*{\fill}
\clearpage

% -----------------------------------------------------------------------------
% Lampiran 1
% -----------------------------------------------------------------------------
\captionlampiran{Rata-rata dan simpangan baku beberapa sifat fisik dan
  kimia tanah dari 78 contoh tanah di Kebun Percobaan Ciheuleut}
\label{lamp:tanah}

\vspace{9.5pt}

{\centering
\fontsize{11pt}{12.65pt}\selectfont
\renewcommand{\arraystretch}{1}%
\begin{tabular}{p{4.52cm}>{\centering\arraybackslash}p{3.66cm}>{\centering\arraybackslash}p{4.55cm}}
  \toprule
  \multicolumn{1}{c}{\textbf{Sifat}}
    & \textbf{Rata-rata}
    & \textbf{Simpangan baku} \\
  \midrule
  Pasir (\%)                        & 47,66 & 23,81 \\
  Lempung (\%)                      & 21,80 & 11,94 \\
  Liat (\%)                         & 30,72 & 18,09 \\
  C-organik (\%)                    & 0,61  & 0,57  \\
  Rapatan isi (mg m$^{-3}$)         & 1,43  & 0,16  \\
  KTK (mek 100 g$^{-1}$ tanah)$^{\mathrm{a}}$
                                     & 18,08 & 17,09 \\
  KAT pada KL (g g$^{-1}$)          & 23,62 & 10,80 \\
  KAT pada TLP (g g$^{-1}$)         & 11,11 & 9,05  \\
  \bottomrule
\end{tabular}\par}

\par\vspace{3pt}\noindent{\footnotesize
  $^{\mathrm{a}}$Banyaknya 70 contoh tanah; KTK: kapasitas tukar kation,
  KAT: kadar air tanah, KL: kapasitas lapang, TLP: titik layu permanen.\par}

\clearpage

% -----------------------------------------------------------------------------
% Lampiran 2
% -----------------------------------------------------------------------------
\captionlampirancentered{Penguasaan pulau berpenghuni di Kepulauan Seribu}
\label{lamp:kepulauan-seribu}

\vspace{6pt}

{\centering
\fontsize{11pt}{13.8pt}\selectfont
\begin{tabular}{@{}>{\centering\arraybackslash}p{1.19cm}@{}p{3.31cm}@{}>{\centering\arraybackslash}p{3.50cm}@{}>{\centering\arraybackslash}p{2cm}@{}>{\centering\arraybackslash}p{2.75cm}@{}}
  \toprule
  \multicolumn{1}{c}{\textbf{No.}}
    & \multicolumn{1}{c}{\textbf{Nama pulau}}
    & \multicolumn{1}{c}{\textbf{Kelurahan}}
    & \multicolumn{1}{c}{\textbf{Luas (ha)}}
    & \multicolumn{1}{c}{\textbf{Penguasaan}} \\
  \specialrule{\lightrulewidth}{0pt}{0pt}
   1 & Pulau Payung      & Pulau Tidung      & 20,86 & Pemprov DKI \\
   2 & Pulau Tidung      &                   & 50,13 & Pemprov DKI \\
   3 & Pulau Lancang     & Pulau Pari        & 15,13 & Pemprov DKI \\
   4 & Pulau Pari        &                   & 41,32 & PT BPA \\
   5 & Pulau Untung Jawa & Pulau Untung Jawa & 40,10 & Pemprov DKI \\
   6 & Pulau Panggang    & Pulau Panggang    &  9,00 & Pemprov DKI \\
   7 & Pulau Pramuka     &                   & 16,00 & Pemprov DKI \\
   8 & Pulau Kelapa      & Pulau Kelapa      & 13,09 & Tidak ada data \\
   9 & Pulau Kelapa Dua  &                   &  1,90 & Tidak ada data \\
  10 & Pulau Harapan     & Pulau Harapan     &  6,70 & Pemprov DKI \\
  11 & Pulau Sebira      &                   &  8,82 & Pemprov DKI \\
  \bottomrule
\end{tabular}\par}

% =============================================================================
% Riwayat Hidup -- Tesis Magister
%
% PPTA 2026: Berisi riwayat hidup, studi, dan kegiatan penulis; maksimum satu
% halaman.
% =============================================================================

\chapter*{RIWAYAT HIDUP}
\addcontentsline{toc}{chapter}{RIWAYAT HIDUP}

Penulis dilahirkan di \TempatLahir{} pada \TanggalLahir{} sebagai anak
ke-\ldots{} dari \ldots{} bersaudara dari pasangan Bapak \ldots{} dan Ibu
\ldots. Pendidikan sarjana (S-1) ditempuh di Program Studi \ldots,
\Fakultas{} \ldots, dan lulus pada tahun \ldots. Pada tahun \ldots, penulis
diterima sebagai mahasiswa program magister (S-2) di \ProgramStudi{} pada
Sekolah Pascasarjana \Universitas{} dan menamatkannya pada tahun \ldots.

Penulis bekerja sebagai \ldots{} di \ldots{} sejak tahun \ldots{} dan
ditempatkan di \ldots{} dan seterusnya (riwayat pekerjaan).

Selama mengikuti program S-2, penulis aktif menjadi \ldots{} (riwayat dan
pengalaman organisasi). Karya ilmiah berjudul \ldots{} telah disajikan pada
seminar/simposium \ldots{} dan/atau dipublikasi di jurnal \ldots{} (riwayat
publikasi). Penulis juga pernah memperoleh atau pernah terpilih sebagai
\ldots{} (riwayat prestasi akademik). Karya-karya ilmiah tersebut merupakan
bagian dari program S-2 penulis.


\end{document}
